\chapter{Detailed Experimental Specifications}
\label{app:experiment_details}

\label{sec:experimental_details}

\noindent\textbf{Structure of Experimental Details}
\label{sec:experiment_details_structure}

This appendix expands each row of Table~\ref{tab:master} into a reproducible specification. To make experiments easy to scan and compare, each experiment section follows a consistent structure:
\begin{itemize}[noitemsep]
    \item \textbf{RQ \& Goal / Hypotheses:} What the experiment is testing and the expected direction of effects.
    \item \textbf{Source Data:} Which datasets/edge tables are used and how ground truth labels are defined (GT Lit vs.\ GT Synth, TP/FP/FN).
    \item \textbf{Simulation Parameters:} What varies vs.\ what is fixed (models, prompts, seeds, corruption rate, judge temperature), aligned with the \textbf{Varies/Fixed} columns in Table~\ref{tab:master}.
    \item \textbf{Simulation/Analysis Procedure:} The operational steps (generation $\rightarrow$ judging $\rightarrow$ correction / Deep Research) and any filtering/deduplication rules.
    \item \textbf{Outputs:} The artefacts produced (Excel/JSON tables, scores, plots) and where they enter downstream analyses.
    \item \textbf{Metrics / Statistical Analysis:} The reported metrics and inference procedure (unit of analysis, tests, and multiple-comparison family).
    \item \textbf{Reproducibility:} The exact scripts/entry points and appendix pointers needed to rerun the experiment end-to-end.
\end{itemize}

When an RQ contains multiple experimental cells (e.g., RQ1a’s judge$\times$GT combinations), we list the shared setup first, then provide a short \textbf{Experiment Details} breakdown for each cell.

%------------------------------------------------------------------------------
% Compact field labels for dense experiment writeups
%------------------------------------------------------------------------------
% In the Experimental Details section, we prefer run-in bold labels over
% heading-style \paragraph blocks to reduce vertical whitespace.
% Experiment-level headers use numbered \subsubsection for clear structure.
\begingroup
\setlength{\leftskip}{1.25em}
\setlength{\parindent}{0pt}
\setlength{\parskip}{2pt}
\makeatletter
\renewcommand{\paragraph}[1]{\par\noindent\textbf{#1}\par\noindent}
\renewcommand{\subparagraph}[1]{\par\noindent\textbf{#1}\par\noindent}
\makeatother

%==============================================================================
\section{Preliminary: Generator Selection}
\label{sec:preliminary_studies}
%==============================================================================

\paragraph{Overview.}
Before running the main experiments, we conducted preliminary studies to select the generator model, prompt, and temperature. These configuration choices are fixed across all subsequent experiments.

%------------------------------------------------------------------------------
\subsubsection{Generator Model Selection}
\label{sec:generator_model_selection}
%------------------------------------------------------------------------------

\paragraph{RQ \& Goal.}
Select a single frontier LLM for all CLD generation experiments, prioritizing edge recovery performance and availability of token-level log probabilities for downstream uncertainty quantification.

\paragraph{Source Data.}
Three validation CLDs (Depressive Symptoms, Social Norms, Emergency Department) with literature-derived ground truth. Node sets are fixed to ground truth variables.
Ground-truth CLD files in Appendix~\ref{sec:data_availability} (Table~\ref{tab:gt_files}).

\paragraph{Simulation Parameters.}
Full configuration in Table~\ref{tab:generator_comparison_config}.

\paragraph{Simulation Procedure.}
For each model--CLD combination, generate a complete \texttt{All Edges} table using the same node set as the literature ground truth. Models are compared purely on edge recovery against GT Lit.
Generation procedure in Appendix~\ref{sec:algorithms} (Algorithm~\ref{alg:cld_generation}).

\paragraph{Simulation Outputs.}
Per-model \texttt{All Edges} tables with relationship types and causal narratives.
Result directory and scripts in Appendix~\ref{sec:generator_comparison} (Table~\ref{tab:generator_comparison_config}).

\paragraph{Unit of Analysis.}
Run-level: each (model, CLD, seed) combination constitutes one observation ($N=27$ total: 3 models $\times$ 3 CLDs $\times$ 3 runs).

\paragraph{Metrics.}
Confusion-matrix counts (TP/FP/FN) computed from \texttt{All Edges} sheets, yielding precision, recall, and F1 (Section~\ref{sec:metrics_definitions}). Aggregate micro-averaged F1 computed by pooling TP/FP/FN across all three CLDs.

\paragraph{Selection Rule.}
Select the generator with highest aggregate (micro-averaged) F1, prioritizing balanced precision/recall. Secondary constraint: token-level logprobs availability for RQ2; among compared models, this is available for GPT-4.1.

\paragraph{Reproducibility.}
Full configuration in Table~\ref{tab:generator_comparison_config}. Execution scripts: Appendix~\ref{sec:code_availability} (Table~\ref{tab:exec_scripts}); analysis scripts: Table~\ref{tab:analysis_scripts}.

\begin{table}[H]
\centering
\caption{Generator Model Comparison Configuration}
\label{tab:generator_comparison_config}
\footnotesize
\begin{tabular}{@{}ll@{}}
\toprule
\textbf{Parameter} & \textbf{Value} \\
\midrule
Models & GPT-4.1, Sonar-Pro, Claude Sonnet 4 \\
CLDs & 3 validation CLDs \\
Generator Prompt & Nitai\_C \\
Temperature & 0.7 \\
Corruption Rate & 0.0 (generation only, no corruption) \\
Judge Edges & False (no judging, compare directly to GT) \\
\midrule
Metric & Micro-averaged Edge F1 across all CLDs \\
Selection Rule & Highest aggregate F1 with logprobs availability \\
\bottomrule
\end{tabular}
\end{table}

\paragraph{Results.}
Table~\ref{tab:generator_model_comparison} (Chapter~\ref{ch:results}).

%------------------------------------------------------------------------------
\subsubsection{Generator Prompt Selection (Nitai\_C)}
\label{sec:generator_prompt_selection}
%------------------------------------------------------------------------------

\paragraph{RQ \& Goal.}
Select a fixed generator prompt variant that maximises edge recovery performance.

\paragraph{Source Data.}
Three validation CLDs (Depressive Symptoms, Social Norms, Emergency Department) with literature-derived ground truth (edge ablation), plus the node-generation ablation dataset described in Appendix~\ref{sec:ablation_node}.
Ground-truth CLD files in Appendix~\ref{sec:data_availability} (Table~\ref{tab:gt_files}).

\paragraph{Simulation Procedure.}
Two ablation studies were conducted:
\begin{itemize}[noitemsep]
    \item \textbf{Edge Generation Ablation:} Nine prompt variants evaluated across three CLDs, testing CoT, task decomposition, RAG/citation requirements, and relationship taxonomy complexity.
    \item \textbf{Node Generation Ablation:} Seven prompt variants evaluated in a 7$\times$3$\times$5 factorial design (105 runs total).
\end{itemize}
Generation procedure in Appendix~\ref{sec:algorithms} (Algorithm~\ref{alg:cld_generation}).

\paragraph{Simulation Outputs.}
Per-variant F1 scores for edge and node generation tasks.
Edge ablation appendix in Appendix~\ref{sec:ablation_edge}; node ablation appendix in Appendix~\ref{sec:ablation_node}.

\paragraph{Unit of Analysis.}
Prompt variant-level: Edge ablation uses 9 prompts $\times$ 3 CLDs = 27 observations; node ablation uses 7 prompts $\times$ 3 CLDs $\times$ 5 runs = 105 observations.

\paragraph{Key Findings.}
\begin{itemize}[noitemsep]
    \item \textbf{Edge generation:} Nitai\_C (dual-level CoT, three relationship types) achieved highest F1 (0.460), outperforming baseline (0.371) by 24\%. CoT improved progressively (Nitai\_A$\rightarrow$B$\rightarrow$C); mandatory citations reduced performance by 25\%; overly restrictive grounding caused complete failure (Cillian\_C: F1=0.00).
    \item \textbf{Node generation:} ANOVA revealed no significant differences ($F=0.81$, $p=0.57$; see Appendix~\ref{app:statistical_procedures_reference}), with F1 ranging 0.617--0.677. Node generation is robust to prompt engineering choices.
\end{itemize}

\paragraph{Reproducibility.}
Exact Nitai\_C prompt in Appendix~\ref{sec:generator_prompt_nitai_c} (Listing~\ref{lst:generator_nitai_c_yaml}). Edge ablation: Appendix~\ref{sec:ablation_edge}. Node ablation: Appendix~\ref{sec:ablation_node}. Execution scripts: Appendix~\ref{sec:code_availability} (Table~\ref{tab:exec_scripts}); analysis scripts: Table~\ref{tab:analysis_scripts}.

\paragraph{Results.}
Edge ablation: Appendix~\ref{sec:ablation_edge}. Node ablation: Appendix~\ref{sec:ablation_node}.

%------------------------------------------------------------------------------
\subsubsection{Generator Temperature Selection}
\label{sec:generator_temp_selection}
%------------------------------------------------------------------------------

\paragraph{RQ \& Goal.}
Validate generator temperature T=0.7 through sensitivity analysis measuring edge recovery performance.

\paragraph{Source Data.}
Three validation CLDs: Social Norms (10 nodes, 12 edges), Depressive Symptoms (14 nodes, 34 edges), Emergency Department (34 nodes, 66 edges).
Ground-truth CLD files in Appendix~\ref{sec:data_availability} (Table~\ref{tab:gt_files}).

\paragraph{Simulation Parameters.}
Full configuration in Table~\ref{tab:temp_sensitivity_config}.

\paragraph{Simulation Procedure.}
For each temperature--CLD--seed combination, generate edges and compute F1 against GT Lit.
Generation procedure in Appendix~\ref{sec:algorithms} (Algorithm~\ref{alg:cld_generation}).

\paragraph{Simulation Outputs.}
Per-run edge F1 scores (45 total).
Temperature sensitivity configuration in Appendix Table~\ref{tab:temp_sensitivity_config}.

\paragraph{Unit of Analysis.}
Run-level: each (temperature, CLD, seed) combination constitutes one observation ($N=45$ total: 5 temperatures $\times$ 3 CLDs $\times$ 3 runs).

\paragraph{Metrics.}
Edge F1 per run; aggregated by temperature level (Section~\ref{sec:metrics_definitions}).

\paragraph{Statistical Analysis.}
One-way ANOVA tested whether temperature significantly affects edge F1 (see Appendix~\ref{app:statistical_procedures_reference}). Results in Table~\ref{tab:temperature_sensitivity}.

\paragraph{Selection Rule.}
Temperature T=0.7 retained as it shows no significant performance difference from other values.

\paragraph{Results.}
Table~\ref{tab:temperature_sensitivity} (Chapter~\ref{ch:results}).

\paragraph{Reproducibility.}
Full configuration in Table~\ref{tab:temp_sensitivity_config}. Execution scripts: Appendix~\ref{sec:code_availability} (Table~\ref{tab:exec_scripts}); analysis scripts: Table~\ref{tab:analysis_scripts}.

\begin{table}[H]
\centering
\caption{Temperature Sensitivity Experiment Configuration}
\label{tab:temp_sensitivity_config}
\footnotesize
\begin{tabular}{@{}ll@{}}
\toprule
\textbf{Parameter} & \textbf{Value} \\
\midrule
CLDs & 3 validation CLDs (Depressive Symptoms, Social Norms, Emergency Department) \\
Generator Model & GPT-4.1 \\
Generator Prompt & Nitai\_C (\texttt{prompts\_Nitai\_C.yaml}) \\
Temperature Values & \{0.0, 0.3, 0.5, 0.7, 1.0\} \\
Runs per Temperature & 3 \\
Seeds & \{10, 20, 30\} \\
Total Runs & 45 (5 temperatures $\times$ 3 seeds $\times$ 3 CLDs) \\
\midrule
\multicolumn{2}{@{}l}{\textit{Metric \& Analysis}} \\
Metric & Edge F1 vs.\ literature ground truth \\
Statistical Test & One-way ANOVA \\
Null Hypothesis & Mean F1 equal across temperatures \\
\bottomrule
\end{tabular}
\end{table}

%==============================================================================
\section{Preliminaries: Generator Performance}
\label{sec:random_baseline_method}
%==============================================================================

\noindent\textbf{RQ \& Goal:}\\
Contextualize generator performance via a structural random baseline---the expected edge-recovery for random graphs with the same $(n,m)$ as each CLD.\\
\textbf{Source Data:}\\
Three validation CLDs with known node and edge counts (Appendix~\ref{sec:data_availability}, Table~\ref{tab:gt_files}).\\
\textbf{Parameters:}\\
Random baseline: 1000 Erd\H{o}s--R\'enyi directed graphs $G(n,m)$ per CLD; LLM generator: 3 runs per CLD (seeds 10, 20, 30). Full configuration in Appendix~\ref{sec:random_baseline} (Table~\ref{tab:random_baseline_config}).\\
\textbf{Unit of Analysis:}\\
Graph-level for random baseline ($N=1000$ per CLD); run-level for LLM ($N=3$ per CLD).\\
\textbf{Simulation Procedure:}\\
For the LLM runs, generate edges via Algorithm~\ref{alg:cld_generation}; for each random graph, sample edges from $G(n,m)$. Compute edge precision/recall/F1 vs.\ the ground truth for both (Table~\ref{tab:random_baseline_config}).\\
\textbf{Outputs:}\\
Random-baseline per-graph F1 distributions (1000 samples per CLD) and per-run LLM F1 values; directory mapping in Appendix~\ref{sec:data_availability} (Table~\ref{tab:experiment_datasets}).\\
\textbf{Metrics/CI:}\\
Precision, recall, F1 (Section~\ref{sec:metrics_definitions}); 95\% CI uses z for random ($n{=}1000$) and t for LLM ($n{=}3$).\\
\textbf{Results:}\\
Table~\ref{tab:random_baseline} and Figure~\ref{fig:random_baseline_final} (Chapter~\ref{ch:results}).\\
\textbf{Reproducibility:}\\
Full configuration in Table~\ref{tab:random_baseline_config}. Execution scripts: Appendix~\ref{sec:code_availability} (Table~\ref{tab:exec_scripts}); analysis scripts: Table~\ref{tab:analysis_scripts}.

\begin{table}[H]
\centering
\caption{Random Baseline Configuration}
\label{tab:random_baseline_config}
\footnotesize
\begin{tabular}{@{}ll@{}}
\toprule
\textbf{Parameter} & \textbf{Value} \\
\midrule
Random Graphs per CLD & 1000 \\
Graph Model & Erd\H{o}s-R\'enyi G(n,m) \\
Node Sets & Same as validation CLDs \\
Edge Counts & Matched to validation CLDs \\
Metric & Edge F1 vs.\ literature ground truth \\
CI Method & z-distribution (n=1000) \\
\bottomrule
\end{tabular}
\end{table}

%------------------------------------------------------------------------------
\section{Preliminary: TruthfulQA Verification}
\label{sec:rq1_verification}
%------------------------------------------------------------------------------

\noindent\textbf{RQ \& Goal:}\\
Verify the correctness-based judge on a general-domain hallucination benchmark before CLD-specific evaluation; strong TruthfulQA performance with weak CLD performance would suggest CLD hallucinations require domain grounding or citation validation beyond correctness judging.\\
\textbf{Hypothesis:}\\
$H_0$: Accuracy $\leq 0.80$; $H_1$: Accuracy $> 0.80$.\\
\textbf{Source Data:}\\
TruthfulQA~\cite{lin2022truthfulqameasuringmodelsmimic} (817 questions). We use the balanced set of 1{,}634 QA pairs provided by TruthfulQA: 817 truthful answers and 817 imitative falsehoods (Appendix~\ref{sec:data_availability}, Table~\ref{tab:validation_datasets}).\\
\textbf{Parameters:}\\
GPT-4.1, temperature 0.0, seed 42, baseline correctness prompt (same structure as RQ1a), max tokens 250. Full configuration in Appendix~\ref{sec:truthqa_verification} (Table~\ref{tab:truthqa_config}).\\

\begin{table}[H]
\centering
\caption{TruthfulQA Verification Configuration}
\label{tab:truthqa_config}
\footnotesize
\begin{tabular}{@{}ll@{}}
\toprule
\textbf{Parameter} & \textbf{Value} \\
\midrule
Dataset & TruthfulQA (817 questions) \\
Samples & 1,634 (817 truthful + 817 hallucinated answers) \\
Judge Model & GPT-4.1 \\
Temperature & 0.0 \\
Seed & 42 \\
Prompt & Baseline correctness \\
Max Tokens & 250 \\
\bottomrule
\end{tabular}
\end{table}
\textbf{Unit of Analysis:}\\
Q\&A pair-level ($N=1634$ total: 817 truthful $+$ 817 hallucinated).\\
\textbf{Simulation Procedure:}\\
The judge labels each answer as \texttt{CORRECT} vs.\ \texttt{INCORRECT} based solely on factual accuracy (Algorithm~\ref{alg:judge_correctness}).\\
\textbf{Outputs:}\\
Per-answer verdicts with reasoning (artefact schema: Appendix~\ref{sec:output_artefacts}, Table~\ref{tab:output_files}).\\
\textbf{Metrics:}\\
Accuracy; Precision/Recall/F1 treating \texttt{Hallucinated} as the positive class (Section~\ref{sec:metrics_rq1a}).\\
\textbf{Reproducibility:}\\
Data in \texttt{final\_runs/RQ1\_verification\_Truth\_QA/}; configuration in Table~\ref{tab:truthqa_config}. Execution and analysis scripts: Appendix~\ref{sec:code_availability} (Tables~\ref{tab:exec_scripts},~\ref{tab:analysis_scripts}).\\
\textbf{Results:}\\
Table~\ref{tab:truthqa_verification} (Chapter~\ref{ch:results}).

%------------------------------------------------------------------------------
\section{RQ1a: LLM-as-a-Judge Hallucination Detection}
\label{sec:rq1a_experiments}
%------------------------------------------------------------------------------

\noindent\textbf{RQ \& Goal:}\\
Test whether LLM judges can detect hallucinated CLD edges: Correctness Judge (Section~\ref{sec:correctness_judge}) and Citation Judge (Section~\ref{sec:citation_judge}) under two ground-truth regimes (GT Lit and GT Synth), yielding four experiment cells.\\
\textbf{Hypotheses:}\\
\textbf{H1 (Correctness):} lower correctness score implies hallucination (FP/FN w.r.t.\ GT).\\
\textbf{H2 (Citation):} lower citation score implies hallucination (FP/FN w.r.t.\ GT).\\
\textbf{Source Data:}\\
Three validation CLDs (Depressive Symptoms, Social Norms, Emergency Department) with literature-derived ground truth; node sets fixed to GT variables (edge-only evaluation; node generation ablation in Appendix~\ref{sec:ablation_node}). Ground-truth files and output directories: Appendix~\ref{sec:data_availability} (Tables~\ref{tab:gt_files},~\ref{tab:experiment_datasets}).\\
\textbf{Parameters:}\\
Correctness judge: 3 prompt variants (\PromptBaseline, \PromptCoT, \PromptMechanistic) $\times$ 3 CLDs $\times$ 3 runs = 27.\\
Citation judge: 4 prompt variants (\PromptBaseline, \PromptCoT, \PromptMechanistic, \PromptMechanisticLit) $\times$ 3 CLDs $\times$ 3 runs = 36.\\
GT Lit seeds 10/20/30 (generation variation); GT Synth seeds 1/2/3 (corruption variation) with 30\% corruption.\\
Citation retrieval in ContentScraper-only mode (Jina AI Reader disabled; Appendix~\ref{alg:citation_fetcher}). Full config in Table~\ref{tab:master}.\\
\textbf{Simulation Procedure:}\\
For each CLD--seed--prompt: generate edges (GPT-4.1 with Nitai\_C); for GT Synth apply Corruptor; run the judge to score each edge; compare scores to GT labels (Algorithms~\ref{alg:cld_generation},~\ref{alg:judge_correctness},~\ref{alg:judge_citation}, and for GT Synth Algorithm~\ref{alg:corruption}).\\
\textbf{Outputs:}\\
Per-edge scores (0.0/0.5/1.0) with verdicts and reasoning; schemas in Appendix~\ref{sec:output_artefacts} (Table~\ref{tab:output_files}).\\
\textbf{Metrics:}\\
AUC-ROC (baseline 0.5), point-biserial correlation $r$, and thresholded classification (score $\leq 0.5$ = hallucination): Precision/Recall/F1 (Section~\ref{sec:metrics_rq1a}).\\
\textbf{Unit of Analysis:}\\
Block-level (CLD$\times$run, $N=9$); edge-level ($N \approx 10{,}000$--$31{,}000$ depending on experiment cell).\\
\textbf{Stats:}\\
Prompt effects under blocked repeated-measures (block = CLD$\times$run, $n{=}9$): Friedman omnibus + paired Wilcoxon with Bonferroni correction (Section~\ref{sec:multiple_comparisons}; Appendix~\ref{app:statistical_procedures_reference}); RM-ANOVA sensitivity decomposition ($\eta^2_p$) in Appendix~\ref{app:sensitivity_analysis}.\\
\textbf{Reproducibility:}\\
Data directories: \texttt{final\_runs/RQ1a\_ground\_truth\_*} and \texttt{final\_runs/RQ1a\_corruption\_*}; see Appendix~\ref{sec:data_availability} (Table~\ref{tab:experiment_datasets}) and Appendix~\ref{sec:code_availability} (Tables~\ref{tab:exec_scripts},~\ref{tab:analysis_scripts}).\\
\textbf{Results:}\\
GT Lit Correctness: Figures~\ref{fig:rq1a_gt_correctness_row1},~\ref{fig:rq1a_gt_correctness_row2}. GT Synth Correctness: Figures~\ref{fig:rq1a_corr_correct_row1},~\ref{fig:rq1a_corr_correct_row2}. GT Lit Citation: Figures~\ref{fig:rq1a_gt_citation_row1},~\ref{fig:rq1a_gt_citation_row2}. GT Synth Citation: Figures~\ref{fig:rq1a_corr_citation_row1},~\ref{fig:rq1a_corr_citation_row2} (Chapter~\ref{ch:results}).

%------------------------------------------------------------------------------
\noindent\textbf{Experiment Details}
%------------------------------------------------------------------------------

\subsubsection{RQ1a: Correctness Judge (Lit)}
\label{sec:rq1a_correctness_lit}

\textbf{RQ \& Goal:}\\
Evaluate Correctness Judge on GT Lit edges (literature-derived ground truth). See shared RQ1a setup (Section~\ref{sec:rq1a_experiments}).\\
\textbf{Source Data:}\\
GT Lit edges from 3 CLDs; ground truth = FP/FN w.r.t.\ expert CLDs. GT files: Appendix Table~\ref{tab:gt_files}.\\
\textbf{Simulation Parameters:}\\
Full configuration in Appendix Table~\ref{tab:rq1a_gt_lit_correctness_config}.\\

\begin{table}[H]
\centering
\caption{RQ1a GT Lit Correctness Judge Configuration}
\label{tab:rq1a_gt_lit_correctness_config}
\footnotesize
\begin{tabular}{@{}ll@{}}
\toprule
\textbf{Parameter} & \textbf{Value} \\
\midrule
Judge Model & GPT-4.1 \\
Temperature & 0.0 \\
Seed & 42 \\
Approach & Correctness-based judging \\
Parallelization & 10 workers \\
Web Search & Disabled \\
\midrule
\multicolumn{2}{@{}l}{\textit{Dataset}} \\
CLDs & Depressive Symptoms, Social Norms, Emergency Department \\
Runs per CLD & 3 (seeds 10, 20, 30) \\
Prompt Variants & 3 (Baseline, CoT, Mechanistic) \\
Total Experiments & 27 (3 CLDs $\times$ 3 runs $\times$ 3 prompts) \\
\bottomrule
\end{tabular}
\end{table}
\textbf{Simulation Procedure:}\\
Algorithm~\ref{alg:cld_generation} (generation) $\rightarrow$ Algorithm~\ref{alg:judge_correctness} (judging).\\
\textbf{Outputs:}\\
Directory: \texttt{final\_runs/RQ1a\_gt\_lit\_correctness/}. Schemas: Appendix Table~\ref{tab:output_files}.\\
\textbf{Unit of Analysis:}\\
CLD$\times$run$\times$prompt ($N=27$: 3 prompts $\times$ 3 CLDs $\times$ 3 runs).\\
\textbf{Metrics / Stats:}\\
See shared RQ1a setup (Section~\ref{sec:metrics_rq1a}); Friedman + Wilcoxon (Appendix~\ref{app:statistical_procedures_reference}).\\
\textbf{Reproducibility:}\\
Scripts: Appendix Tables~\ref{tab:exec_scripts},~\ref{tab:analysis_scripts}.\\
\textbf{Results:}\\
Figures~\ref{fig:rq1a_gt_correctness_row1},~\ref{fig:rq1a_gt_correctness_row2} (Chapter~\ref{ch:results}).

\subsubsection{RQ1a: Correctness Judge (Synth)}
\label{sec:rq1a_correctness_synth}

\textbf{RQ \& Goal:}\\
Evaluate Correctness Judge on GT Synth edges (30\% corruption). See shared RQ1a setup (Section~\ref{sec:rq1a_experiments}).\\
\textbf{Source Data:}\\
GT Synth edges from 3 CLDs; ground truth = corruption status (corrupted = hallucination). Dataset mapping: Appendix~\ref{sec:data_availability} (Table~\ref{tab:experiment_datasets}).\\
\textbf{Simulation Parameters:}\\
Full configuration in Appendix Table~\ref{tab:rq1a_gt_synth_correctness_config}.\\

\begin{table}[H]
\centering
\caption{RQ1a GT Synth Correctness Judge Configuration}
\label{tab:rq1a_gt_synth_correctness_config}
\footnotesize
\begin{tabular}{@{}ll@{}}
\toprule
\textbf{Parameter} & \textbf{Value} \\
\midrule
Judge Model & GPT-4.1 \\
Temperature & 0.0 \\
Seed & 42 \\
Approach & Correctness-based judging \\
Parallelization & 10 workers \\
Web Search & Disabled \\
\midrule
\multicolumn{2}{@{}l}{\textit{Corruption}} \\
Corruption Rate & 30\% \\
Corruption Types & Motivation (spurious), Structural (new edges, polarity flips) \\
\midrule
\multicolumn{2}{@{}l}{\textit{Dataset}} \\
CLDs & Depressive Symptoms, Social Norms, Emergency Department \\
Runs per CLD & 3 (seeds 1, 2, 3) \\
Prompt Variants & 3 (Baseline, CoT, Mechanistic) \\
Total Experiments & 27 (3 CLDs $\times$ 3 runs $\times$ 3 prompts) \\
\bottomrule
\end{tabular}
\end{table}
\textbf{Simulation Procedure:}\\
Algorithm~\ref{alg:cld_generation} (generation) $\rightarrow$ Algorithm~\ref{alg:corruption} (corruption) $\rightarrow$ Algorithm~\ref{alg:judge_correctness} (judging).\\
\textbf{Outputs:}\\
Directory: \texttt{final\_runs/RQ1a\_gt\_synth\_correctness/}. Schemas: Appendix Table~\ref{tab:output_files}.\\
\textbf{Unit of Analysis:}\\
CLD$\times$run$\times$prompt ($N=27$: 3 prompts $\times$ 3 CLDs $\times$ 3 runs).\\
\textbf{Metrics / Stats:}\\
See shared RQ1a setup (Section~\ref{sec:metrics_rq1a}); Friedman + Wilcoxon (Appendix~\ref{app:statistical_procedures_reference}).\\
\textbf{Reproducibility:}\\
Scripts: Appendix Tables~\ref{tab:exec_scripts},~\ref{tab:analysis_scripts}.\\
\textbf{Results:}\\
Figures~\ref{fig:rq1a_corr_correct_row1},~\ref{fig:rq1a_corr_correct_row2} (Chapter~\ref{ch:results}).

\subsubsection{RQ1a: Citation Judge (Lit)}
\label{sec:rq1a_citation_lit}

\textbf{RQ \& Goal:}\\
Evaluate Citation Judge on GT Lit edges (literature retrieval). See shared RQ1a setup (Section~\ref{sec:rq1a_experiments}).\\
\textbf{Source Data:}\\
GT Lit edges from 3 CLDs; ground truth = FP/FN w.r.t.\ expert CLDs. GT files: Appendix Table~\ref{tab:gt_files}.\\
\textbf{Simulation Parameters:}\\
Full configuration in Appendix Table~\ref{tab:rq1a_gt_lit_citation_config}.\\

\begin{table}[H]
\centering
\caption{RQ1a GT Lit Citation Judge Configuration}
\label{tab:rq1a_gt_lit_citation_config}
\footnotesize
\begin{tabular}{@{}ll@{}}
\toprule
\textbf{Parameter} & \textbf{Value} \\
\midrule
Judge Model & GPT-4.1 \\
Temperature & 0.0 \\
Seed & 42 \\
Approach & Citation-based judging \\
Parallelization & 10 workers \\
Web Search & Enabled \\
Citation Fetcher & ContentScraper (Jina AI Reader disabled) \\
\midrule
\multicolumn{2}{@{}l}{\textit{Dataset}} \\
CLDs & Depressive Symptoms, Social Norms, Emergency Department \\
Runs per CLD & 3 (seeds 10, 20, 30) \\
Prompt Variants & 4 (Baseline, CoT, Mechanistic, Mechanistic-Lit) \\
Total Experiments & 36 (3 CLDs $\times$ 3 runs $\times$ 4 prompts) \\
\bottomrule
\end{tabular}
\end{table}
\textbf{Simulation Procedure:}\\
Algorithm~\ref{alg:cld_generation} (generation) $\rightarrow$ Algorithm~\ref{alg:citation_fetcher} (retrieval) $\rightarrow$ Algorithm~\ref{alg:judge_citation} (judging).\\
\textbf{Outputs:}\\
Directory: \texttt{final\_runs/RQ1a\_gt\_lit\_citation/}. Schemas: Appendix Table~\ref{tab:output_files}.\\
\textbf{Unit of Analysis:}\\
CLD$\times$run$\times$prompt ($N=36$: 4 prompts $\times$ 3 CLDs $\times$ 3 runs).\\
\textbf{Metrics / Stats:}\\
See shared RQ1a setup (Section~\ref{sec:metrics_rq1a}); Friedman + Wilcoxon (Appendix~\ref{app:statistical_procedures_reference}).\\
\textbf{Reproducibility:}\\
Scripts: Appendix Tables~\ref{tab:exec_scripts},~\ref{tab:analysis_scripts}.\\
\textbf{Results:}\\
Figures~\ref{fig:rq1a_gt_citation_row1},~\ref{fig:rq1a_gt_citation_row2} (Chapter~\ref{ch:results}).

\subsubsection{RQ1a: Citation Judge (Synth)}
\label{sec:rq1a_citation_synth}

\textbf{RQ \& Goal:}\\
Evaluate Citation Judge on GT Synth edges (30\% corruption, GPT-5-mini). See shared RQ1a setup (Section~\ref{sec:rq1a_experiments}).\\
\textbf{Source Data:}\\
GT Synth edges from 3 CLDs; ground truth = corruption status (corrupted = hallucination). Dataset mapping: Appendix~\ref{sec:data_availability} (Table~\ref{tab:experiment_datasets}).\\
\textbf{Simulation Parameters:}\\
Full configuration in Appendix Table~\ref{tab:gt_synth_citation_config}.\\

\begin{table}[H]
\centering
\caption{GT Synth Citation Judge Configuration (GPT-5-mini)}
\label{tab:gt_synth_citation_config}
\footnotesize
\begin{tabular}{@{}ll@{}}
\toprule
\textbf{Parameter} & \textbf{Value} \\
\midrule
Judge Model & GPT-5-mini \\
Temperature & 0.0 \\
Top-p & 1.0 \\
Approach & Citation-based judging \\
Parallelization & 10 workers \\
\midrule
\multicolumn{2}{@{}l}{\textit{Corruption}} \\
Corruption Rate & 30\% \\
Corruption Types & Motivation (spurious explanations), Structural (new edges, polarity flips) \\
\midrule
\multicolumn{2}{@{}l}{\textit{Dataset}} \\
CLDs & Depressive Symptoms, Social Norms, Emergency Department \\
Runs per CLD & 3 (seeds 1, 2, 3) \\
Prompt Variants & 4 (Baseline, CoT, Mechanistic, Mechanistic-Lit) \\
Total Experiments & 36 (3 CLDs $\times$ 3 runs $\times$ 4 prompts) \\
\bottomrule
\end{tabular}
\end{table}
\textbf{Simulation Procedure:}\\
Algorithm~\ref{alg:cld_generation} (generation) $\rightarrow$ Algorithm~\ref{alg:corruption} (corruption) $\rightarrow$ Algorithm~\ref{alg:citation_fetcher} (retrieval) $\rightarrow$ Algorithm~\ref{alg:judge_citation} (judging).\\
\textbf{Outputs:}\\
Directory: \texttt{final\_runs/RQ1a\_gt\_synth\_citation/}. Schemas: Appendix Table~\ref{tab:output_files}.\\
\textbf{Unit of Analysis:}\\
CLD$\times$run$\times$prompt ($N=36$: 4 prompts $\times$ 3 CLDs $\times$ 3 runs).\\
\textbf{Metrics / Stats:}\\
See shared RQ1a setup (Section~\ref{sec:metrics_rq1a}); Friedman + Wilcoxon (Appendix~\ref{app:statistical_procedures_reference}).\\
\textbf{Reproducibility:}\\
Scripts: Appendix Tables~\ref{tab:exec_scripts},~\ref{tab:analysis_scripts}.\\
\textbf{Results:}\\
Figures~\ref{fig:rq1a_corr_citation_row1},~\ref{fig:rq1a_corr_citation_row2} (Chapter~\ref{ch:results}).

%------------------------------------------------------------------------------
\section{Sensitivity: Prompt Sensitivity Analysis}
\label{sec:rq1a_sensitivity}
%------------------------------------------------------------------------------

\paragraph{RQ \& Goal.}
Quantify prompt engineering impact on judge performance using variance-based sensitivity analysis inspired by Sobol indices~\cite{sobol1993sensitivity}.

\paragraph{Source Data.}
Prompt-level and block-level results from the four RQ1a judge experiments (summarised in Table~\ref{tab:master}).
Sensitivity output directory in Appendix~\ref{sec:data_availability} (Table~\ref{tab:experiment_datasets}).

\paragraph{Simulation Parameters.}
Full configuration in Table~\ref{tab:prompt_sensitivity_config}.

\begin{table}[H]
\centering
\caption{Prompt Sensitivity Analysis Configuration}
\label{tab:prompt_sensitivity_config}
\footnotesize
\begin{tabular}{@{}ll@{}}
\toprule
\textbf{Parameter} & \textbf{Value} \\
\midrule
Embedding Model & MPNet (all-mpnet-base-v2, 768-dim) \\
Distance Metric & Cosine distance from baseline prompt \\
Prompt Variants & Judge: Baseline, CoT, Mechanistic (+ Mechanistic-Lit for citation); Corrector: Baseline, CoT, Mechanistic \\
Judge Experiments & 4 (Corruption Correctness/Citation, GT Correctness/Citation) \\
Corrector Experiments & 2 (Synthetic Corruption, Ground Truth) \\
Statistical Test & Friedman (blocked by CLD$\times$run), post-hoc paired Wilcoxon (Bonferroni) \\
\bottomrule
\end{tabular}
\end{table}

\paragraph{Simulation Procedure.}
Repeated-measures ANOVA with CLD$\times$run as subjects and prompt as within-subject factor.
Full RM-ANOVA tables and diagnostics in Appendix~\ref{app:sensitivity_analysis}.

\paragraph{Simulation Outputs.}
RM-ANOVA effect sizes ($\eta^2_p$), assumption diagnostics, and sensitivity figures for each judge$\times$GT setting.
Output artefacts in Appendix~\ref{sec:output_artefacts} (Table~\ref{tab:output_files}); sensitivity output directory in Appendix~\ref{sec:data_availability} (Table~\ref{tab:experiment_datasets}).

\paragraph{Metrics.}
Following Sobol's variance-based sensitivity framework, for a categorical factor $A\in\{1,\dots,K\}$ (here: prompt choice) the first-order index is:
\begin{equation}
S_A := \frac{\mathrm{Var}(\mathbb{E}[Y \mid A])}{\mathrm{Var}(Y)}.
\end{equation}
For categorical $A$, this reduces to $S_A=\eta^2$,
i.e., the population one-way ANOVA effect size (correlation ratio) for a categorical factor~\cite{saltelli2010variance}. In the main
experiments, we report $\eta^2_p$ (partial eta-squared) because our design is blocked by CLD$\times$run (repeated
measures), so the sensitivity measure is computed after controlling for these subject-level differences.
Partial eta-squared: $\eta^2_p = \frac{SS_{\text{prompt}}}{SS_{\text{prompt}} + SS_{\text{error}}}$. Measures proportion of within-subject variance in F1 explained by prompt choice. Interpretation (Cohen): $<$0.01 negligible, 0.01--0.06 small, 0.06--0.14 medium, $\geq$0.14 large.

\paragraph{Unit of Analysis.}
Block-level (CLD$\times$run, $N=9$) for RM-ANOVA variance decomposition; prompt as within-subject factor.

\paragraph{Statistical Analysis.}
RM-ANOVA assumption diagnostics: Shapiro--Wilk (residual normality), Mauchly (sphericity), Greenhouse--Geisser correction when violated, and Bonferroni correction (see Appendix~\ref{app:statistical_procedures_reference}). Full tables in Appendix~\ref{app:sensitivity_analysis}.

\paragraph{Reproducibility.}
Data: \texttt{final\_runs/supp\_prompt\_sensitivity/}. Dataset mapping in Appendix~\ref{sec:data_availability} (Table~\ref{tab:experiment_datasets}). Execution and analysis scripts: Appendix~\ref{sec:code_availability} (Tables~\ref{tab:exec_scripts},~\ref{tab:analysis_scripts}).

\paragraph{Results.}
Appendix~\ref{app:sensitivity_analysis} (Figure~\ref{fig:sensitivity_analysis_app}).

%------------------------------------------------------------------------------
\section{Validation: Uleman's Provider Study}
\label{sec:rq1a_validation_studies}
%------------------------------------------------------------------------------

\paragraph{Overview.}
Four validation studies assess judge reliability: (1) search provider comparison, (2) human-LLM agreement, (3) physics CLD sanity check, and (4) retrieval success analysis.

\paragraph{Source Data.}
Validation datasets (Uleman's CLD, physics CLDs, human validation set) and their output directories.
Dataset locations in Appendix~\ref{sec:data_availability} (Table~\ref{tab:validation_datasets}).

\paragraph{Simulation Procedure.}
Apply the relevant judge pipeline(s) (correctness/citation) to validation edges; citation judging includes retrieval and fetching.
Algorithms in Appendix~\ref{sec:algorithms} (Algorithms~\ref{alg:judge_correctness},~\ref{alg:judge_citation},~\ref{alg:citation_fetcher}).

\paragraph{Simulation Outputs.}
Per-edge verdicts/scores (and retrieval diagnostics for citation judging) stored in the corresponding validation output directories.
Output artefacts in Appendix~\ref{sec:output_artefacts} (Table~\ref{tab:output_files}); validation directories in Appendix~\ref{sec:data_availability} (Table~\ref{tab:validation_datasets}).

\paragraph{Reproducibility.}
Dataset specifications: Appendix~\ref{sec:data_availability} (Tables~\ref{tab:validation_datasets},~\ref{tab:physics_cld_dataset}). Execution and analysis scripts: Appendix~\ref{sec:code_availability} (Tables~\ref{tab:exec_scripts},~\ref{tab:analysis_scripts}).

%------------------------------------------------------------------------------
\subsection{Search Provider Comparison (Uleman's CLD)}
\label{sec:ulemans_validation}
%------------------------------------------------------------------------------

\paragraph{RQ \& Goal.}
Validate citation-based judge on an independently expert-validated dataset. Determine whether low GT Lit scores are due to search infrastructure limitations.

\paragraph{Hypothesis.}
If low scores are infrastructure-driven, we expect significant variation across providers/fetchers. Consistently low scores would confirm a fundamental validation gap between expert reasoning and literature-based verification.

\paragraph{Source Data.}
Uleman's CLD: 150-edge causal loop diagram for Alzheimer's disease, constructed through Group Model Building (GMB) with 17 domain experts~\cite{uleman2024triangulation}. Unlike primary CLDs (extracted from literature), represents direct expert consensus with bibliographic references. Dataset: Appendix~\ref{sec:data_availability} (Table~\ref{tab:validation_datasets}).

\paragraph{Simulation Parameters.}
Full configuration in Table~\ref{tab:ulemans_config}.

\begin{table}[H]
\centering
\caption{Uleman's CLD Validation Configuration}
\label{tab:ulemans_config}
\footnotesize
\begin{tabular}{@{}ll@{}}
\toprule
\textbf{Parameter} & \textbf{Value} \\
\midrule
Dataset & Uleman's Alzheimer's CLD (150 edges) \\
Search Providers & OpenAlex, Brave Search (with/without whitelist), Perplexity, PubMed \\
Citation Fetchers & Jina AI Reader, ContentScraper (BeautifulSoup-based, PDF-capable) \\
Configurations & 4 providers $\times$ 2 fetchers = 8 \\
Total Judgements & 1,200 (8 configs $\times$ 150 edges) \\
\midrule
\multicolumn{2}{@{}l}{\textit{Judge Settings}} \\
Judge Model & GPT-4.1 \\
Temperature & 0.0 \\
Top-p & 1.0 \\
Prompt & Baseline citation prompt \\
\bottomrule
\end{tabular}
\end{table}

\paragraph{Simulation Procedure.}
For each provider--fetcher configuration, retrieve citations and apply the Citation Judge to all 150 edges (Algorithm~\ref{alg:citation_fetcher} $\rightarrow$ Algorithm~\ref{alg:judge_citation}).

\paragraph{Outputs.}
Per-edge scores and provider comparison summaries; output artefacts in Appendix~\ref{sec:output_artefacts} (Table~\ref{tab:output_files}).

\paragraph{Metrics.}
Average score (mean 0--1), coverage (successful fetch \%), score difference (ContentScraper $-$ Jina AI); see Section~\ref{sec:metrics_validation}.

\paragraph{Unit of Analysis.}
Edge-level ($N=150$); crossed with 4 providers $\times$ 2 fetchers = 1,200 total judgements.

\paragraph{Reproducibility.}
Data: \texttt{final\_runs/validation\_RQ1a\_ulemans\_external/}. Configuration: Table~\ref{tab:ulemans_config}. Execution and analysis scripts: Appendix~\ref{sec:code_availability} (Tables~\ref{tab:exec_scripts},~\ref{tab:analysis_scripts}).

\paragraph{Results.}
Table~\ref{tab:ulemans_search_providers} (Chapter~\ref{ch:results}).

%------------------------------------------------------------------------------
\subsection{Human--LLM Agreement}
\label{sec:human_llm_agreement}
%------------------------------------------------------------------------------

\paragraph{RQ \& Goal.}
Validate that LLM citation judges produce reliable verdicts by measuring inter-rater agreement with human experts.

\paragraph{Hypothesis.}
Expect substantial agreement ($\kappa > 0.60$), performance comparable to human baseline, and consistency across CLDs/generators.

\paragraph{Source Data.}
72 causal edges across 3 validation files (2 CLDs, 2 generators): 30 edges Depressive Symptoms/GPT-4.1, 30 edges Depressive Symptoms/Claude 4 Sonnet, 12 edges Social Norms/GPT-4.1. Dataset: Appendix~\ref{sec:data_availability} (Table~\ref{tab:validation_datasets}).

\paragraph{Simulation Parameters.}
Human validator (domain expert) and LLM judge (GPT-4.1, T=0, baseline prompt) each provide categorical verdicts: Not supported (0), Partially supported (0.5), Fully supported (1). Judge settings follow RQ1a Citation Judge configuration. Full configuration in Table~\ref{tab:human_validation_config}.

\begin{table}[H]
\centering
\caption{Human-LLM Agreement Validation Configuration}
\label{tab:human_validation_config}
\footnotesize
\begin{tabular}{@{}ll@{}}
\toprule
\textbf{Parameter} & \textbf{Value} \\
\midrule
Sample Size & 72 edges across 3 validation files \\
CLDs & Depressive Symptoms, Social Norms \\
Generators & GPT-4.1, Claude 4 Sonnet \\
Human Validator & Domain expert (categorical verdicts) \\
LLM Judge & Citation-based (GPT-4.1, T=0.0, baseline) \\
\midrule
\multicolumn{2}{@{}l}{\textit{Metrics}} \\
Agreement & Cohen's $\kappa$, Agreement Rate, Pearson $r$ \\
Interpretation & \cite{landis1977measurement} guidelines \\
\bottomrule
\end{tabular}
\end{table}

\paragraph{Simulation Procedure.}
For each edge, retrieve citations and apply the Citation Judge (Algorithm~\ref{alg:citation_fetcher} $\rightarrow$ Algorithm~\ref{alg:judge_citation}); human validator independently annotates the same edges.

\paragraph{Outputs.}
Per-edge human and LLM verdicts with agreement statistics; output artefacts in Appendix~\ref{sec:output_artefacts} (Table~\ref{tab:output_files}).

\paragraph{Metrics.}
Metrics defined in Section~\ref{sec:metrics_validation}:
\begin{itemize}[noitemsep]
    \item \textbf{Cohen's $\kappa$:} Chance-corrected agreement (Appendix~\ref{app:statistical_procedures_reference}). Interpretation~\citep{landis1977measurement}: $<$0.20 slight, 0.21--0.40 fair, 0.41--0.60 moderate, 0.61--0.80 substantial, $>$0.80 almost perfect.
    \item \textbf{Agreement Rate:} Exact match percentage
    \item \textbf{Pearson $r$:} Score correlation (Appendix~\ref{app:statistical_procedures_reference})
\end{itemize}
Baseline: Human inter-rater reliability from prior GMB studies ($\kappa \approx 0.60$, 80\% agreement).

\paragraph{Unit of Analysis.}
Edge-level ($N=72$ total across 3 validation files: 30 + 30 + 12 edges).

\paragraph{Reproducibility.}
Data: \texttt{final\_runs/RQ1\_human\_validation\_citation\_judge/}. Configuration: Table~\ref{tab:human_validation_config}. Execution and analysis scripts: Appendix~\ref{sec:code_availability} (Tables~\ref{tab:exec_scripts},~\ref{tab:analysis_scripts}).

\paragraph{Results.}
Tables~\ref{tab:human_validation},~\ref{tab:human_validation_confusion} (Chapter~\ref{ch:results}).

%------------------------------------------------------------------------------
\subsection{Physics Sanity Check}
\label{sec:physics_validation}
%------------------------------------------------------------------------------

\paragraph{RQ \& Goal.}
Verify that judges correctly identify valid causal relationships when unambiguous (sanity check under curated inputs).

\paragraph{Source Data.}
Four physics-based CLDs derived from fundamental laws (31 edges, 26 variables):
\begin{enumerate}[noitemsep]
    \item Thermostat Heating System: Newton's Law of Cooling~\cite{newton1701scala}, First Law of Thermodynamics~\cite{clausius1850heat}
    \item Predator-Prey Dynamics: Lotka-Volterra equations~\cite{lotka1925elements,volterra1928fluctuations}
    \item RC Circuit: Ohm's Law~\cite{ohm1827galvanische}, Kirchhoff's Laws~\cite{kirchhoff1845durchgang}
    \item Water Tank: Pascal's Law~\cite{pascal1663equilibre}, Torricelli's Law~\cite{torricelli1644opera}
\end{enumerate}
CLD structure (variables, edges, polarities) and edge rationales were authored in an interactive session with Claude 4.5 Opus~\cite{anthropic2025claude45opus} by translating these laws into CLD form, then verified against source texts. Per-CLD breakdown: Appendix Table~\ref{tab:physics_cld_dataset}.

\paragraph{Metrics.} We report \emph{approval} rates computed from discrete judge verdicts (see Section~\ref{sec:metrics_validation}). For the Correctness Judge, approval is the fraction of edges with verdict \texttt{CORRECT} (score 1.0). For the Citation Judge, approval is the fraction of edges with aggregate verdict \texttt{FULLY\_SUPPORTED} or \texttt{PARTIALLY\_SUPPORTED} (score $\ge 0.5$); \texttt{NO\_CITATION} indicates retrieval failure. We report approval across all edges and conditional approval restricted to edges with at least one retrieved citation ($C_e\ge1$).

\paragraph{Simulation Parameters.}
Full configuration in Table~\ref{tab:physics_cld_config}. Judge settings follow RQ1a configuration (GPT-4.1, T=0.0).

\begin{table}[H]
\centering
\caption{Physics CLD Validation Configuration}
\label{tab:physics_cld_config}
\footnotesize
\begin{tabular}{@{}ll@{}}
\toprule
\textbf{Parameter} & \textbf{Value} \\
\midrule
CLDs & Thermostat (8 edges), Predator-Prey (10), RC Circuit (7), Water Tank (6) \\
Total Edges & 31 \\
Judge Model & GPT-4.1 \\
Temperature & 0.0 \\
\midrule
\multicolumn{2}{@{}l}{\textit{Physics Principles}} \\
Thermostat & Newton's Law of Cooling, First Law of Thermodynamics \\
Predator-Prey & Lotka--Volterra population dynamics \\
RC Circuit & Ohm's Law, Kirchhoff's Laws \\
Water Tank & Pascal's Law, Torricelli's Law \\
\bottomrule
\end{tabular}
\end{table}

\paragraph{Simulation Procedure.}
Both Correctness Judge and Citation Judge applied to curated physics edges (Algorithms~\ref{alg:judge_correctness},~\ref{alg:judge_citation}).

\paragraph{Outputs.}
Per-edge verdicts and scores; output artefacts in Appendix~\ref{sec:output_artefacts} (Table~\ref{tab:output_files}).

\paragraph{Unit of Analysis.}
Edge-level ($N=31$ edges across 4 physics CLDs).

\paragraph{Reproducibility.}
Data: \texttt{final\_runs/validation\_RQ1a\_physics\_clds/}. Configuration: Table~\ref{tab:physics_cld_config}. Execution and analysis scripts: Appendix~\ref{sec:code_availability} (Tables~\ref{tab:exec_scripts},~\ref{tab:analysis_scripts}).

\paragraph{Results.}
Table~\ref{tab:physics-comparison-enhanced} (Chapter~\ref{ch:results}).

%------------------------------------------------------------------------------
\subsection{Retrieval Success Analysis}
\label{sec:retrieval_success}
%------------------------------------------------------------------------------

\paragraph{RQ \& Goal.}
Investigate whether low citation scores are driven by retrieval failures (missing documents) or evidence characteristics (weak/hedged claims).

\paragraph{Source Data.}
\begin{enumerate}[noitemsep]
    \item Expert-provided citations: Physics CLDs, Uleman's CLD (150 edges)
    \item Automated web search: Three main validation CLDs via generate-then-search pipeline
\end{enumerate}
For Uleman's and Physics CLDs, CLD structure and causal narratives were authored with Claude 4.5 Opus~\cite{anthropic2025claude45opus}. Dataset mapping: Appendix~\ref{sec:data_availability} (Table~\ref{tab:validation_datasets}).

\paragraph{Simulation Procedure.}
Apply the citation retrieval pipeline to each dataset (Algorithm~\ref{alg:citation_fetcher}); compute fetch success rate and compare against judge approval rates.

\paragraph{Metrics.}
Retrieval success rate = percentage of edges with successfully fetched content (see Section~\ref{sec:metrics_validation}). Discrepancies between retrieval success and judge approval distinguish infrastructure failures from validation failures. Fetch success criteria: Appendix~\ref{alg:citation_fetcher}.

\paragraph{Unit of Analysis.}
Edge-level (diagnostic analysis comparing retrieval success vs.\ approval rates across datasets).

\paragraph{Reproducibility.}
Analysis integrated with Uleman's and Physics CLD validation studies. Execution and analysis scripts: Appendix~\ref{sec:code_availability} (Tables~\ref{tab:exec_scripts},~\ref{tab:analysis_scripts}).

\paragraph{Results.}
Table~\ref{tab:retrieval-comparison} (Chapter~\ref{ch:results}).

%------------------------------------------------------------------------------
\section{RQ1b: LLM-as-a-Corrector Evaluation}
\label{sec:rq1b}
\label{sec:rq1b_experiments}
%------------------------------------------------------------------------------

\noindent\textbf{RQ \& Goal:}\\
Evaluate whether the Corrector (Section~\ref{sec:corrector}) can repair judge-flagged hallucinations.\\
\textbf{Hypothesis:}\\
H3: correction improves edge quality for flagged edges ($\Delta$F1 $> 0$).\\
\textbf{Source Data:}\\
Judge-flagged edges from RQ1a (GT Synth and GT Lit); output directories in Appendix~\ref{sec:data_availability} (Table~\ref{tab:experiment_datasets}).\\
\textbf{Parameters:}\\
3 corrector prompts (Baseline, CoT, Mechanistic); design = 3 prompts $\times$ 3 CLDs $\times$ 3 runs = 27 runs per GT type; evaluation judge = baseline Correctness Judge; scope = correctness-based correction only (citation-corrector omitted due to retrieval cost). Full configuration in Table~\ref{tab:master}.\\
\textbf{Simulation Procedure:}\\
For each flagged edge, apply the Corrector (revise relationship type and/or narrative), re-judge the corrected edge, and compare pre/post metrics against ground truth (Algorithms~\ref{alg:corrector} and~\ref{alg:judge_correctness}).\\
\textbf{Outputs:}\\
Corrected edges plus re-judge scores; artefact schema in Appendix~\ref{sec:output_artefacts} (Table~\ref{tab:output_files}).\\
\textbf{Metrics:}\\
$\Delta$F1/$\Delta$Precision/$\Delta$Recall (pre$\rightarrow$post vs.\ GT) and $\Delta$JudgeScore = mean(corrected) $-$ mean(original); see Section~\ref{sec:metrics_rq1b}.\\
\textbf{Unit of Analysis:}\\
Block-level (CLD$\times$run, $N=9$) for statistical inference; edge-level for individual correction evaluation.\\
\textbf{Stats:}\\
Efficacy via one-sample Wilcoxon signed-rank on block-level $\Delta$F1 (blocks = CLD$\times$run, $n{=}9$) testing median$(\Delta\mathrm{F1})=0$; prompt effects via Friedman omnibus + paired Wilcoxon with Bonferroni correction (Section~\ref{sec:multiple_comparisons}; Appendix~\ref{app:statistical_procedures_reference}); RM-ANOVA sensitivity in Appendix~\ref{app:sensitivity_analysis}.\\
\textbf{Reproducibility:}\\
Data in \texttt{final\_runs/RQ1b\_corrector\_experiment\_*/}; scripts in Appendix~\ref{sec:code_availability} (Tables~\ref{tab:exec_scripts},~\ref{tab:analysis_scripts}); analysis script \texttt{analyze\_corrector\_ablation.py}.\\
\textbf{Results:}\\
GT Synth: Tables~\ref{tab:rq1b_synthetic_correctness},~\ref{tab:rq1b_prompt_ablation_synthetic}. GT Lit: Tables~\ref{tab:rq1b_groundtruth_correctness},~\ref{tab:rq1b_prompt_ablation_groundtruth} (Chapter~\ref{ch:results}).

%------------------------------------------------------------------------------
\noindent\textbf{Experiment Details}
%------------------------------------------------------------------------------

\subsubsection{RQ1b: Correctness Corrector (Synth)}
\label{sec:rq1b_corrector_synth}

\textbf{RQ \& Goal:}\\
Evaluate Corrector on GT Synth flagged edges. See shared RQ1b setup (Section~\ref{sec:rq1b_experiments}).\\
\textbf{Source Data:}\\
Judge-flagged edges from RQ1a GT Synth experiments (corrupted CLDs). Dataset mapping: Appendix~\ref{sec:data_availability} (Table~\ref{tab:experiment_datasets}).\\
\textbf{Simulation Parameters:}\\
Full configuration in Appendix Table~\ref{tab:rq1b_corrector_synth_config}.\\

\begin{table}[H]
\centering
\caption{RQ1b Corrector (GT Synth) Configuration}
\label{tab:rq1b_corrector_synth_config}
\footnotesize
\begin{tabular}{@{}ll@{}}
\toprule
\textbf{Parameter} & \textbf{Value} \\
\midrule
\multicolumn{2}{@{}l}{\textit{Corrector}} \\
Corrector Model & GPT-4.1 \\
Temperature & 0.0 \\
Corrector Prompts & 3 (Baseline, CoT, Mechanistic) \\
\midrule
\multicolumn{2}{@{}l}{\textit{Evaluation Judge}} \\
Judge Model & GPT-4.1 \\
Judge Prompt & Baseline Correctness \\
Temperature & 0.0 \\
\midrule
\multicolumn{2}{@{}l}{\textit{Dataset}} \\
Input & RQ1a GT Synth flagged edges \\
CLDs & Depressive Symptoms, Social Norms, Emergency Department \\
Runs per CLD & 3 (seeds 1, 2, 3) \\
Total Experiments & 27 (3 CLDs $\times$ 3 runs $\times$ 3 prompts) \\
\bottomrule
\end{tabular}
\end{table}
\textbf{Simulation Procedure:}\\
Algorithm~\ref{alg:corrector} (correction) $\rightarrow$ Algorithm~\ref{alg:judge_correctness} (re-judging).\\
\textbf{Outputs:}\\
Directory: \texttt{final\_runs/RQ1b\_corrector\_experiment\_corruption\_*/}. Schemas: Appendix Table~\ref{tab:output_files}.\\
\textbf{Unit of Analysis:}\\
Block-level (CLD$\times$run, $N=9$); see shared RQ1b setup (Section~\ref{sec:rq1b_experiments}).\\
\textbf{Metrics / Stats:}\\
$\Delta$F1 pre$\rightarrow$post (Section~\ref{sec:metrics_rq1b}); Wilcoxon signed-rank (Appendix~\ref{app:statistical_procedures_reference}).\\
\textbf{Reproducibility:}\\
Scripts: Appendix Tables~\ref{tab:exec_scripts},~\ref{tab:analysis_scripts}.\\
\textbf{Results:}\\
Tables~\ref{tab:rq1b_synthetic_correctness},~\ref{tab:rq1b_prompt_ablation_synthetic} (Chapter~\ref{ch:results}).

\subsubsection{RQ1b: Correctness Corrector (Lit)}
\label{sec:rq1b_corrector_lit}

\textbf{RQ \& Goal:}\\
Evaluate Corrector on GT Lit flagged edges. See shared RQ1b setup (Section~\ref{sec:rq1b_experiments}).\\
\textbf{Source Data:}\\
Judge-flagged edges from RQ1a GT Lit experiments. GT files: Appendix Table~\ref{tab:gt_files}.\\
\textbf{Simulation Parameters:}\\
Full configuration in Appendix Table~\ref{tab:rq1b_corrector_lit_config}.\\

\begin{table}[H]
\centering
\caption{RQ1b Corrector (GT Lit) Configuration}
\label{tab:rq1b_corrector_lit_config}
\footnotesize
\begin{tabular}{@{}ll@{}}
\toprule
\textbf{Parameter} & \textbf{Value} \\
\midrule
\multicolumn{2}{@{}l}{\textit{Corrector}} \\
Corrector Model & GPT-4.1 \\
Temperature & 0.0 \\
Corrector Prompts & 3 (Baseline, CoT, Mechanistic) \\
\midrule
\multicolumn{2}{@{}l}{\textit{Evaluation Judge}} \\
Judge Model & GPT-4.1 \\
Judge Prompt & Baseline Correctness \\
Temperature & 0.0 \\
\midrule
\multicolumn{2}{@{}l}{\textit{Dataset}} \\
Input & RQ1a GT Lit flagged edges \\
CLDs & Depressive Symptoms, Social Norms, Emergency Department \\
Runs per CLD & 3 (seeds 10, 20, 30) \\
Total Experiments & 27 (3 CLDs $\times$ 3 runs $\times$ 3 prompts) \\
\bottomrule
\end{tabular}
\end{table}
\textbf{Simulation Procedure:}\\
Algorithm~\ref{alg:corrector} (correction) $\rightarrow$ Algorithm~\ref{alg:judge_correctness} (re-judging).\\
\textbf{Outputs:}\\
Directory: \texttt{final\_runs/RQ1b\_corrector\_experiment\_ground\_truth\_*/}. Schemas: Appendix Table~\ref{tab:output_files}.\\
\textbf{Unit of Analysis:}\\
Block-level (CLD$\times$run, $N=9$); see shared RQ1b setup (Section~\ref{sec:rq1b_experiments}).\\
\textbf{Metrics / Stats:}\\
$\Delta$F1 pre$\rightarrow$post (Section~\ref{sec:metrics_rq1b}); Wilcoxon signed-rank (Appendix~\ref{app:statistical_procedures_reference}).\\
\textbf{Reproducibility:}\\
Scripts: Appendix Tables~\ref{tab:exec_scripts},~\ref{tab:analysis_scripts}.\\
\textbf{Results:}\\
Tables~\ref{tab:rq1b_groundtruth_correctness},~\ref{tab:rq1b_prompt_ablation_groundtruth} (Chapter~\ref{ch:results}).

%------------------------------------------------------------------------------
\section{RQ2: Corpus-Independent Hallucination Detection}
\label{sec:rq2_experiments}
%------------------------------------------------------------------------------

\noindent\textbf{RQ \& Goal:}\\
Test whether UQ metrics captured during generation (Section~\ref{sec:uq_metrics_definition}) can predict hallucinations without external context.\\
\textbf{Hypotheses:}\\
\textbf{H4 (UQ):} higher perplexity / lower $P_{\min}$ / higher window entropy $\Rightarrow$ hallucination (block-level AUC $>0.5$).\\
\textbf{H5 (CosSim):} lower narrative--citation cosine similarity $\Rightarrow$ hallucination (block-level AUC $>0.5$).\\
\textbf{Source Data:}\\
63 deduplicated GT Lit files (36 citation, 27 correctness), 31{,}704 edges, hallucination rate 15.2\%. Hallucination label: (FP or FN) w.r.t.\ GT Lit. \textit{Metric availability:} Logprob-derived metrics (perplexity, min prob, max window entropy) available for all 63 files; cosine similarity available only for citation-judging runs (35 files, 16{,}492 edges) because it requires retrieved citation text. Ensemble phases use a 16{,}507-edge subset (35 citation files) where all four metrics are available. Outputs and GT files: Appendix~\ref{sec:data_availability} (Tables~\ref{tab:experiment_datasets},~\ref{tab:gt_files}).\\
\textbf{UQ Metrics + Parameters:}\\
Computed from generator logprobs (Algorithms~\ref{alg:uq_metrics},~\ref{alg:classification}). Perplexity uses capped NLL averaging (token cap 20.0; average cap 10.0, bounding PPL to $\le e^{10}\approx22{,}026$); $P_{\min}$ is the minimum token probability; $H_{\text{win}}$ is max sliding-window entropy (window size 5 over top-$k$ entropies with $k{=}5$); CosSim is the max narrative--citation-chunk embedding similarity for citation experiments (embedding model all-mpnet-base-v2 from \cite{muennighoff2023mteb}). Full shared configuration in Table~\ref{tab:rq2_shared_config}.

\begin{table}[H]
\centering
\caption{RQ2 Shared Experimental Configuration}
\label{tab:rq2_shared_config}
\footnotesize
\begin{tabular}{@{}ll@{}}
\toprule
\textbf{Parameter} & \textbf{Value} \\
\midrule
CI Metrics & Gen Perplexity, Gen Min Prob, Gen Max Window Entropy, Gen Cosine Similarity \\
Block Definition & CLD $\times$ run ($N=9$: 3 CLDs $\times$ 3 runs) \\
Hallucination Label & FP or FN w.r.t.\ GT Lit \\
Scaler & StandardScaler \\
Random State & 42 \\
\bottomrule
\end{tabular}
\end{table}
\textbf{Outputs:}\\
Per-edge metric values, per-file AUCs, and classifier predictions; schemas in Appendix~\ref{sec:output_artefacts} (Table~\ref{tab:output_files}).\\
\textbf{Metrics:}\\
AUC-ROC, point-biserial correlation $r$, and classifier CV/test AUC (Section~\ref{sec:metrics_rq2}). Because hallucinations are the minority class (15.2\%), AUC is used as a threshold-free separation metric; thresholded Precision/Recall depend on the chosen operating point and class prevalence.\\
\textbf{Unit of Analysis:}\\
Block-level (CLD$\times$run, $N=9$) for statistical inference; edge-level ($N=31{,}704$) for descriptive analysis.\\
\textbf{Stats:}\\
Block-level aggregation (CLD$\times$run, $N{=}9$) as the unit of inference. \textit{Single-metric analysis:} one-sample \textbf{Wilcoxon signed-rank} tests on $(\mathrm{AUC}-0.5)$ with Bonferroni correction across 4 metrics ($m{=}4$, $\alpha_{\text{adj}}{=}0.0125$); correlation $r$ aggregated via Fisher $z$-transform across blocks with separate Bonferroni family. \textit{Edge-level:} Mann--Whitney U tests reported as \textbf{Significant Files (\%)}---the percentage of files where edge-level $p<0.05$; this is a \textit{descriptive} consistency measure only (not used for confirmatory inference due to within-file dependence). \textit{File inclusion:} files excluded if $<$10 edges or $<$2 per class (hallucination/correct). \textit{Uncertainty:} 95\% CIs via $t$-distribution over blocks ($df{=}8$). Full procedures: Appendix~\ref{app:statistical_procedures_reference}.\\
\textbf{Reproducibility:}\\
Data in \texttt{final\_runs/RQ2\_uq\_hallucination\_detection/}; dataset mapping in Appendix~\ref{sec:data_availability} (Table~\ref{tab:experiment_datasets}). Execution and analysis scripts: Appendix~\ref{sec:code_availability} (Tables~\ref{tab:exec_scripts},~\ref{tab:analysis_scripts}).\\
\textbf{Results:}\\
Figures~\ref{fig:rq2_distributions_by_cld},~\ref{fig:rq2_rfe}; Tables~\ref{tab:rq2_single_metrics},~\ref{tab:rq2_normality_tests},~\ref{tab:rq2_ensemble_performance} (Chapter~\ref{ch:results}).

%------------------------------------------------------------------------------
\noindent\textbf{Experiment Details}
%------------------------------------------------------------------------------

\subsubsection{RQ2: UQ (Logprobs)}
\label{sec:rq2_uq_logprobs}

\textbf{RQ \& Goal:}\\
Test whether logprob-derived UQ metrics can predict hallucinations. See shared RQ2 setup (Section~\ref{sec:rq2_experiments}).\\
\textbf{Source Data:}\\
RQ1a GT Lit experiments (27 correctness + 36 citation files, 31,704 edges); hallucination = FP/FN. Dataset mapping: Appendix~\ref{sec:data_availability} (Table~\ref{tab:experiment_datasets}).\\
\textbf{Simulation Parameters:}\\
Full configuration in Appendix Table~\ref{tab:rq2_uq_logprobs_config}.\\

\begin{table}[H]
\centering
\caption{RQ2 UQ (Logprobs) Configuration}
\label{tab:rq2_uq_logprobs_config}
\footnotesize
\begin{tabular}{@{}ll@{}}
\toprule
\textbf{Parameter} & \textbf{Value} \\
\midrule
\multicolumn{2}{@{}l}{\textit{UQ Metrics}} \\
Perplexity (PPL) & Capped NLL averaging (token cap 20.0, avg cap 10.0) \\
$P_{\min}$ & Minimum token probability \\
$H_{\text{win}}$ & Max sliding-window entropy (window 5, top-$k$=5) \\
\midrule
\multicolumn{2}{@{}l}{\textit{Dataset}} \\
Source & RQ1a GT Lit experiments (27 correctness files) \\
Total Edges & 31,704 (from 63 deduplicated files) \\
Hallucination Rate & 15.2\% \\
\midrule
\multicolumn{2}{@{}l}{\textit{Analysis}} \\
Block Definition & CLD $\times$ run ($N=9$ blocks) \\
Statistical Test & One-sample $t$-test vs.\ 0.5 \\
\bottomrule
\end{tabular}
\end{table}
\textbf{Simulation Procedure:}\\
Algorithm~\ref{alg:uq_metrics} (metric extraction from logprobs).\\
\textbf{Outputs:}\\
Directory: \texttt{final\_runs/RQ2\_uq\_hallucination\_detection/}. Schemas: Appendix Table~\ref{tab:output_files}.\\
\textbf{Unit of Analysis:}\\
Block-level (CLD$\times$run, $N=9$); see shared RQ2 setup (Section~\ref{sec:rq2_experiments}).\\
\textbf{Metrics / Stats:}\\
Block-level AUC (Section~\ref{sec:metrics_rq2}); one-sample \textbf{Wilcoxon signed-rank} test on $(\mathrm{AUC}-0.5)$ (Appendix~\ref{app:statistical_procedures_reference}).\\
\textbf{Reproducibility:}\\
Scripts: Appendix Tables~\ref{tab:exec_scripts},~\ref{tab:analysis_scripts}.\\
\textbf{Results:}\\
Figure~\ref{fig:rq2_distributions_by_cld}; Tables~\ref{tab:rq2_single_metrics},~\ref{tab:rq2_normality_tests} (Chapter~\ref{ch:results}).

\subsubsection{RQ2: CosSim (Embeddings)}
\label{sec:rq2_cossim_embeddings}

\textbf{RQ \& Goal:}\\
Test whether narrative--citation cosine similarity can predict hallucinations. See shared RQ2 setup (Section~\ref{sec:rq2_experiments}).\\
\textbf{Source Data:}\\
RQ1a GT Lit Citation experiments only (36 files); ensemble subset = 16,507 edges. Dataset mapping: Appendix~\ref{sec:data_availability} (Table~\ref{tab:experiment_datasets}).\\
\textbf{Simulation Parameters:}\\
Full configuration in Appendix Table~\ref{tab:rq2_cossim_config}.\\

\begin{table}[H]
\centering
\caption{RQ2 Cosine Similarity (Embeddings) Configuration}
\label{tab:rq2_cossim_config}
\footnotesize
\begin{tabular}{@{}ll@{}}
\toprule
\textbf{Parameter} & \textbf{Value} \\
\midrule
\multicolumn{2}{@{}l}{\textit{Embedding Model}} \\
Model & all-mpnet-base-v2 \\
Source & sentence-transformers (Hugging Face) \\
Similarity & Max narrative--citation-chunk cosine similarity \\
\midrule
\multicolumn{2}{@{}l}{\textit{Dataset}} \\
Source & RQ1a GT Lit Citation experiments only (36 files) \\
Ensemble Subset & 16,507 edges (35 files with all 4 metrics) \\
Hallucination Rate & 15.2\% \\
\midrule
\multicolumn{2}{@{}l}{\textit{Analysis}} \\
Block Definition & CLD $\times$ run ($N=9$ blocks) \\
Statistical Test & One-sample $t$-test vs.\ 0.5 \\
\bottomrule
\end{tabular}
\end{table}
\textbf{Simulation Procedure:}\\
Compute narrative--citation-chunk embeddings and cosine similarity via Algorithm~\ref{alg:uq_metrics}. Benchmark: Appendix~\ref{sec:embedding_benchmark}.\\
\textbf{Outputs:}\\
Directory: \texttt{final\_runs/RQ2\_uq\_hallucination\_detection/}. Schemas: Appendix Table~\ref{tab:output_files}.\\
\textbf{Unit of Analysis:}\\
Block-level (CLD$\times$run, $N=9$); see shared RQ2 setup (Section~\ref{sec:rq2_experiments}).\\
\textbf{Metrics / Stats:}\\
Block-level AUC (Section~\ref{sec:metrics_rq2}); one-sample \textbf{Wilcoxon signed-rank} test on $(\mathrm{AUC}-0.5)$ (Appendix~\ref{app:statistical_procedures_reference}).\\
\textbf{Reproducibility:}\\
Scripts: Appendix Tables~\ref{tab:exec_scripts},~\ref{tab:analysis_scripts}.\\
\textbf{Results:}\\
Figure~\ref{fig:rq2_distributions_by_cld}; Table~\ref{tab:rq2_single_metrics} (Chapter~\ref{ch:results}).

%------------------------------------------------------------------------------
\subsubsection{RQ2: Analysis Pipeline (6 Phases)}
\label{sec:rq2_analysis}
%------------------------------------------------------------------------------

\noindent\textbf{Simulation Procedure (6 phases):}\\
\textbf{Phases 1--3 (single-metric analysis):}\\
\textit{Data loading:} Load 63 deduplicated experiment files from inventory (36 citation, 27 correctness); exclude corruption experiments.\\
\textit{Hallucination labeling:} \texttt{is\_hallucination} = True if Classification $\in$ \{FP, FN\} w.r.t.\ GT Lit.\\
\textit{Data cleaning:} Remove rows with missing or infinite CI metric values (dropna on feature columns + label).\\
\textit{Per-file analysis:} Compute AUC-ROC for each file (hallucination vs.\ correct); require $\geq$10 edges and $\geq$2 per class.\\
\textit{Block aggregation:} Group files by (CLD $\times$ run) to form $N{=}9$ blocks; compute mean AUC per block.\\
\textit{Statistical test:} One-sample Wilcoxon signed-rank test on $(\mathrm{AUC}-0.5)$ across blocks (Bonferroni: $m{=}4$, $\alpha_{\text{adj}}{=}0.0125$; Appendix~\ref{app:statistical_procedures_reference}).\\
Full configuration in Table~\ref{tab:rq2_phases123_config}.

\begin{table}[H]
\centering
\caption{RQ2 Phases 1--3 (Single-Metric Analysis) Configuration}
\label{tab:rq2_phases123_config}
\footnotesize
\begin{tabular}{@{}ll@{}}
\toprule
\textbf{Parameter} & \textbf{Value} \\
\midrule
\multicolumn{2}{@{}l}{\textit{Data Loading}} \\
Input Files & 63 deduplicated experiment files (36 citation, 27 correctness) \\
Filter & Exclude corruption experiments \\
\midrule
\multicolumn{2}{@{}l}{\textit{Hallucination Labeling}} \\
Label & is\_hallucination = True if Classification $\in$ \{FP, FN\} \\
Ground Truth & GT Lit (literature-based correctness) \\
\midrule
\multicolumn{2}{@{}l}{\textit{Data Cleaning}} \\
Missing Values & Remove rows with NaN in CI metric columns \\
Infinite Values & Replace $\pm\infty$ with NaN, then drop \\
Min Edges/File & $\geq$10 edges required \\
Min Class Size & $\geq$2 per class (hallucination / correct) required \\
\midrule
\multicolumn{2}{@{}l}{\textit{Block Aggregation}} \\
Block Key & (CLD, run) $\Rightarrow$ $N=9$ blocks \\
Aggregation & Mean AUC per block \\
\midrule
\multicolumn{2}{@{}l}{\textit{Statistical Testing}} \\
Test & One-sample Wilcoxon signed-rank on (AUC $-$ 0.5) \\
Alternative & Two-sided \\
Bonferroni Correction & $m=4$ tests, $\alpha_{\text{adj}}=0.0125$ \\
\bottomrule
\end{tabular}
\end{table}
\textbf{Phase 4 (RFE):} Recursive Feature Elimination~\cite{pedregosa2011sklearn} with \texttt{GroupKFold(n\_splits=3)} over blocks; classifiers: LogReg, RF, GradBoost.\\
\textbf{Phase 5 (ensemble):} block-level train/test split (67\%/33\%) via \texttt{GroupShuffleSplit} and model selection on training blocks via \texttt{GroupKFold}; four classifiers including a neural network.\\
\textbf{Classifier families:} brief algorithm notes for LR/RF/GB/MLP appear in Appendix~\ref{sec:rq2_classifier_notes}.\\
\textbf{Phase 6 (LOCO):} leave-one-CLD-out evaluation for cross-domain generalisation. Full formal specification in Appendix~\ref{alg:classification}.\\
\textbf{Simulation Parameters:}\\
Full pipeline configuration in Table~\ref{tab:rq2_analysis_pipeline_config}.

\begin{table}[H]
\centering
\caption{RQ2 Analysis Pipeline Configuration}
\label{tab:rq2_analysis_pipeline_config}
\footnotesize
\begin{tabular}{@{}lll@{}}
\toprule
\textbf{Phase} & \textbf{Parameter} & \textbf{Value} \\
\midrule
\multirow{3}{*}{Phase 4 (RFE)} & Classifiers & LogReg, RandomForest, GradientBoosting \\
 & Cross-validation & GroupKFold ($n\_splits=3$) \\
 & Feature range & $k \in \{1,2,3,4\}$ \\
\midrule
\multirow{4}{*}{Phase 5 (Ensemble)} & Classifiers & LogReg, RF, GB, MLP \\
 & Train/Test Split & 67\%/33\% (GroupShuffleSplit) \\
 & CV on training & 3-fold GroupKFold \\
 & Class weighting & balanced (LogReg, RF) \\
\midrule
\multirow{2}{*}{Phase 6 (LOCO)} & Protocol & Leave-one-CLD-out \\
 & Scaling & StandardScaler (fit on training CLDs) \\
\bottomrule
\end{tabular}
\end{table}
\textbf{Reproducible splitting protocol (Phases 4--6):}\\
Grouping key is \texttt{block\_id=cld\_run} (3 CLDs $\times$ 3 runs $\Rightarrow N{=}9$); rows are deterministically sorted by \((\texttt{block\_id}, \texttt{source\_file})\) prior to any group split; scaling uses \texttt{StandardScaler} fit on training blocks only.\\
Phase 4 holds out 3 blocks and trains on 6; evaluates \(k \in \{1,2,3,4\}\) and selects the \(k\) maximizing validation AUC per fold.\\
Phase 5 holds out 3 blocks for testing and uses the remaining 6 for training; CV on training uses 3 folds (2 blocks held out per fold). For F1@$t^*$ reporting, the threshold $t^*$ is selected by maximizing out-of-fold F1 on training blocks only; this fixed $t^*$ is then frozen and applied to both the Phase 5 test set and all Phase 6 evaluations.\\
Phase 6 holds out an entire CLD as test and trains on the other two CLDs (scaling fit on training CLDs only). 95\% CIs are computed via the $t$-distribution with $df=n-1$ over folds ($n=3$ for both phases).\\
\textbf{Unit of Analysis:}\\
Block-level (CLD$\times$run, $N=9$) for Phases 1--5; CLD-level ($N=3$) for Phase 6 (LOCO).\\
\textbf{Limitations:}\\
UQ metrics are not applicable for GT Synth (corruption is post-generation); cosine similarity is only defined for Citation Judge runs.\\
\textbf{Reproducibility:}\\
Data in \texttt{final\_runs/RQ2\_uq\_hallucination\_detection/}; dataset mapping in Appendix~\ref{sec:data_availability} (Table~\ref{tab:experiment_datasets}). Execution and analysis scripts: Appendix~\ref{sec:code_availability} (Tables~\ref{tab:exec_scripts},~\ref{tab:analysis_scripts}).\\
\textbf{Results:}\\
Figures~\ref{fig:rq2_distributions_by_cld},~\ref{fig:rq2_rfe}; Table~\ref{tab:rq2_ensemble_performance}; Section~\ref{sec:rq2_ensemble} (Chapter~\ref{ch:results}).
    
%------------------------------------------------------------------------------
\section{RQ3: Deep Research Literature Validation}
\label{sec:rq3_experiments}
%------------------------------------------------------------------------------

\begin{center}
\fbox{\parbox{0.95\textwidth}{
\textbf{Methodological Note:} The original RQ3 design proposed \textit{smart test-time compute allocation}---using UQ metrics to selectively deploy the LLM-as-a-corrector. This was predicated on: (1) RQ2 UQ metrics generalizing across CLDs, and (2) RQ1b correctors effectively repairing hallucinations. Neither condition was met (see Discussion, Section~\ref{sec:discussion_rq3}). We therefore pivoted to \textbf{Deep Research validation} (system documented in Section~\ref{sec:deep_research_system}).

This section presents Deep Research as an \textbf{exploratory pilot study}. Given time constraints from the late pivot, the experimental design is less comprehensive than RQ1--RQ2: single-run (no replication), limited parameter ablation, and no end-to-end pipeline integration.
}}
\end{center}

% (spacing trimmed for density)

\subsubsection{Deep Research}
\label{sec:rq3_deep_research}

\paragraph{RQ \& Goal.}
Evaluate whether Deep Research (Section~\ref{sec:deep_research_system}) can distinguish valid from hallucinated causal claims by finding \emph{direct causal} evidence in the scientific literature.

\paragraph{Hypothesis.}
H6: high detection for TP, low detection for FP, and a TP--FP gap $>$20 percentage points.

\paragraph{Source Data.}
285 edges from RQ1a GT Lit experiments (TP=44, FP=178, FN=63). Output directory: Appendix~\ref{sec:data_availability} (Table~\ref{tab:experiment_datasets}).

\paragraph{Parameters.}
Single run per edge ($\sim\$0.87$/edge) with \texttt{max\_iterations=1}, \texttt{max\_sources=5}, \texttt{MAX\_SEARCHES=5}, \texttt{time\_limit=900s}.

\paragraph{Simulation Procedure.}
Generate search queries from the causal claim, run parallel Brave searches, score evidence, apply strict direct-causality criteria, and synthesize a verdict (Algorithm~\ref{alg:deep_research}).

\paragraph{Outputs.}
Per-edge verdict, confidence (1--10), \texttt{direct\_causal\_found} flag, and source URLs; schemas in Appendix~\ref{sec:output_artefacts} (Table~\ref{tab:output_files}).

\paragraph{Metrics.}
Detection rate of \texttt{direct\_causal\_found}; TP--FP gap $\Delta>20$pp; mean confidence by class (Section~\ref{sec:metrics_rq3}).

\paragraph{Stats.}
Pairwise Fisher exact tests (Bonferroni: $m=3$, $\alpha_{\text{adj}}=0.017$); effect size Cohen's $h$ (Appendix~\ref{app:statistical_procedures_reference}).

\paragraph{Unit of Analysis.}
Edge-level ($N=285$ edges: TP=44, FP=178, FN=63).

\paragraph{Reproducibility.}
Full configuration in Table~\ref{tab:rq3_config}. Execution and analysis scripts: Appendix~\ref{sec:code_availability} (Tables~\ref{tab:exec_scripts},~\ref{tab:analysis_scripts}).

\begin{table}[H]
\centering
\caption{RQ3 Deep Research Configuration}
\label{tab:rq3_config}
\footnotesize
\begin{tabular}{@{}ll@{}}
\toprule
\textbf{Parameter} & \textbf{Value} \\
\midrule
Heavy Model & GPT-5 (hypothesis evaluation, judgement, verdict) \\
Light Model & GPT-5-mini (iterative search, evidence scoring) \\
Temperature & 1.0 (both models) \\
Top-p & 1.0 (both models) \\
Search API & Brave Search with academic domain filtering \\
Max Iterations & 1 (system supports up to 3) \\
Queries per Iteration & 3--4 (generated by hypothesis\_evaluator) \\
Max Searches per Query & 5 \\
Articles Selected & 2--3 (for full-text fetch) \\
Time Limit & 900 seconds (15 minutes) per edge \\
Batch Size & 5 edges processed concurrently \\
\midrule
\multicolumn{2}{@{}l}{\textit{Dataset}} \\
Total Edges & 285 (44 TP, 178 FP, 63 FN) \\
CLDs & Depressive Symptoms (84), Social Norms (17), Emergency Department (184) \\
\midrule
\multicolumn{2}{@{}l}{\textit{Human Validation}} \\
Sample Size & 50 edges (29.6\% of edges with direct\_causal\_found=True) \\
Sampling & Stratified by CLD $\times$ classification (seed=42) \\
Thresholds Tested & \{0.5, 0.6, 0.7, 0.8, 0.9\} \\
\bottomrule
\end{tabular}
\end{table}

\paragraph{Results.}
Figure~\ref{fig:rq3_combined}, Table~\ref{tab:scaling_summary} (Chapter~\ref{ch:results}).

%------------------------------------------------------------------------------
\subsubsection{Computational Cost}
\label{sec:rq3_computational_cost}
%------------------------------------------------------------------------------

\noindent\textbf{Cost summary (reported in Results).}\\
See Results Section~\ref{sec:rq3_cost_results} (Table~\ref{tab:scaling_summary}) for the Deep Research computational cost breakdown and the overall cost comparison (\$0.87/edge; $\sim$64$\times$ citation judging; $\sim$544$\times$ correctness judging).

%------------------------------------------------------------------------------
\subsubsection{Limitations}
\label{sec:rq3_limitations}
%------------------------------------------------------------------------------

\begin{itemize}[noitemsep]
    \item \textbf{No replication:} Single run per edge prevents uncertainty quantification
    \item \textbf{Evidence applicability:} Retrieved causal evidence may be only partially aligned to specific edge claims
    \item \textbf{Full-text access:} Paywalled articles limit evidence quality
    \item \textbf{Search API dependency:} Results depend on Brave Search indexing
\end{itemize}

%------------------------------------------------------------------------------
\subsubsection{Exploratory Analysis: Ground Truth Enrichment}
\label{sec:gt_enrichment_method}
%------------------------------------------------------------------------------

\paragraph{RQ \& Goal.}
Explore whether DR results could enrich expert-constructed GT by validating ``hallucinations'' that may represent valid causal relationships experts omitted.

\paragraph{Source Data.}
RQ3 Deep Research outputs (285 edges). Dataset mapping: Appendix~\ref{sec:data_availability} (Table~\ref{tab:experiment_datasets}).

\paragraph{Extrapolation Procedure.}
Two enrichment scenarios:
\begin{itemize}[noitemsep]
    \item \textbf{Scenario 1 (Enriched GT):} Promote DR-validated FPs to TP in ground truth; recompute Precision/Recall/F1.
    \item \textbf{Scenario 2 (LLM+DR):} Additionally add DR-validated FNs to system output.
\end{itemize}

\paragraph{Simulation Parameters.}
Full configuration in Table~\ref{tab:gt_enrichment_config}.

\begin{table}[H]
\centering
\caption{Ground Truth Enrichment Configuration}
\label{tab:gt_enrichment_config}
\footnotesize
\begin{tabular}{@{}ll@{}}
\toprule
\textbf{Parameter} & \textbf{Value} \\
\midrule
\multicolumn{2}{@{}l}{\textit{Scenario 1 (Enriched GT)}} \\
FP Promotion & DR-validated FPs $\rightarrow$ TP in ground truth \\
Metrics Recomputed & Precision, Recall, F1 \\
\midrule
\multicolumn{2}{@{}l}{\textit{Scenario 2 (LLM+DR)}} \\
FP Promotion & Same as Scenario 1 \\
FN Addition & DR-validated FNs added to system output \\
Metrics Recomputed & Precision, Recall, F1 \\
\midrule
\multicolumn{2}{@{}l}{\textit{Dataset}} \\
Source & RQ3 Deep Research outputs \\
Total Edges & 285 (44 TP, 178 FP, 63 FN) \\
\bottomrule
\end{tabular}
\end{table}

\paragraph{Caveats.}
Domain variation (28\% Depressive Symptoms vs 76\% Emergency Department FP promotion), evidence specificity concerns, circularity risk (LLM validating LLM), no expert review. Results exploratory.

\paragraph{Results.}
Section~\ref{sec:gt_enrichment_results}, Table~\ref{tab:enriched_gt} (Chapter~\ref{ch:results}).

\paragraph{Reproducibility.}
\texttt{python final\_runs/RQ3\_deep\_research\_validation/compute\_enriched\_gt\_metrics.py}. Analysis scripts: Appendix~\ref{sec:code_availability} (Table~\ref{tab:analysis_scripts}).

%------------------------------------------------------------------------------
\subsubsection{Human Validation}
\label{sec:rq3_human_validation}
%------------------------------------------------------------------------------

\paragraph{RQ \& Goal.}
Validate DR verdicts via human review to establish applicability thresholds and enable extrapolation.

\paragraph{Source Data.}
50 edges (29.6\% of 169 eligible) via proportional stratified sampling across CLD $\times$ classification. Seed 42 for reproducibility. Dataset: Appendix~\ref{sec:data_availability} (Table~\ref{tab:experiment_datasets}).

\paragraph{Annotation Procedure.}
Human reviewer assigns categorical verdict (\texttt{causal}, \texttt{different\_vars}, \texttt{invalid}) and continuous applicability score $\in[0,1]$ using the applicability rubric in Appendix Table~\ref{tab:human-validation-rubric-rq3}: 1.0 = exact match; 0.8--0.9 = different operationalisations; $\approx$0.7 = partial claim; $\approx$0.6 = proxy constructs; $\leq$0.5 = weak/non-causal evidence.

\paragraph{Simulation Parameters.}
Full configuration in Table~\ref{tab:rq3_human_validation_config}.

\begin{table}[H]
\centering
\caption{RQ3 Human Validation Configuration}
\label{tab:rq3_human_validation_config}
\footnotesize
\begin{tabular}{@{}ll@{}}
\toprule
\textbf{Parameter} & \textbf{Value} \\
\midrule
Sample Size & 50 edges (29.6\% of 169 eligible) \\
Random Seed & 42 \\
Sampling Method & Proportional stratified (CLD $\times$ classification) \\
Filter & direct\_causal\_found = True \\
Verdict Categories & causal, different\_vars, invalid \\
Applicability Score & Continuous $\in [0, 1]$ \\
\bottomrule
\end{tabular}
\end{table}

\paragraph{Results.}
Figure~\ref{fig:rq3_human_validation_elbow}, Table~\ref{tab:rq3_human_validation_tradeoff}, Table~\ref{tab:rq3_verdict_distribution} (Chapter~\ref{ch:results}).

\paragraph{Reproducibility.}
\texttt{python final\_runs/RQ3\_deep\_research\_validation/scripts/sample\_edges\_for\_validation.py}. Analysis scripts: Appendix~\ref{sec:code_availability} (Table~\ref{tab:analysis_scripts}).

\paragraph{Inter-rater reliability check.}
To assess rubric reliability, a second rater independently annotated an overlapping subset of 13 edges using the same rubric. Agreement for the categorical verdict was quantified using Cohen's $\kappa$ (chance-corrected; interpretation per~\cite{landis1977measurement} as defined in Section~\ref{sec:human_llm_agreement}; Appendix~\ref{app:statistical_procedures_reference}). Agreement for the continuous applicability score was quantified using the Intraclass Correlation Coefficient, ICC(2,1) (two-way random effects, absolute agreement, single measures; Appendix~\ref{app:statistical_procedures_reference})~\citep{koo2016guideline}. The ICC is defined as:
\[
\text{ICC} = \frac{\sigma^2_{\text{between}}}{\sigma^2_{\text{between}} + \sigma^2_{\text{within}}}
\]
where $\sigma^2_{\text{between}}$ is the variance between subjects (edges) and $\sigma^2_{\text{within}}$ is the variance within subjects (across raters). Interpretation thresholds~\citep{koo2016guideline}: $<$0.50 poor, 0.50--0.75 moderate, 0.75--0.90 good, $>$0.90 excellent.\\
\textbf{Reproducibility:} \texttt{python final\_runs/RQ3\_deep\_research\_validation/validation/calculate\_inter\_rater\_reliability.py}

%------------------------------------------------------------------------------
\subsubsection{Threshold Selection and Extrapolation}
\label{sec:rq3_threshold_selection}
%------------------------------------------------------------------------------

\paragraph{RQ \& Goal.}
Determine applicability threshold for filtering DR verdicts and extrapolate enrichment potential to full dataset.

\paragraph{Source Data.}
Human-validated sample (50 edges from RQ3 Human Validation). Dataset mapping: Appendix~\ref{sec:data_availability} (Table~\ref{tab:experiment_datasets}).

\paragraph{Extrapolation Procedure.}
Sensitivity analysis across $t \in \{0.5, 0.6, 0.7, 0.8, 0.9\}$. For each threshold, compute coverage (pass rate) and correct-causal rate with Wilson 95\% CI (Appendix~\ref{app:statistical_procedures_reference}). Select threshold via gradient-based elbow rule on coverage--correctness curve.

\paragraph{Simulation Parameters.}
Full configuration in Table~\ref{tab:rq3_threshold_config}.

\begin{table}[H]
\centering
\caption{RQ3 Threshold Selection Configuration}
\label{tab:rq3_threshold_config}
\footnotesize
\begin{tabular}{@{}ll@{}}
\toprule
\textbf{Parameter} & \textbf{Value} \\
\midrule
Thresholds Tested & $t \in \{0.5, 0.6, 0.7, 0.8, 0.9\}$ \\
Selection Method & Gradient-based elbow rule \\
CI Method & Wilson 95\% CI \\
\midrule
\multicolumn{2}{@{}l}{\textit{Metrics at Each Threshold}} \\
Coverage & $x(t) = \Pr(\text{applicability} \geq t)$ \\
Correct-Causal Rate & $y(t) = \Pr(\texttt{causal} \mid \text{applicability} \geq t)$ \\
\midrule
\multicolumn{2}{@{}l}{\textit{Extrapolation}} \\
Sample Size & 50 edges (29.6\% of 169 eligible) \\
Full Dataset & 285 edges \\
\bottomrule
\end{tabular}
\end{table}

\paragraph{Metrics.}
Metrics defined in Section~\ref{sec:metrics_rq3}:
\begin{itemize}[noitemsep]
    \item \textbf{Coverage:} $x(t) = \Pr(\text{applicability} \geq t)$
    \item \textbf{Correct-causal rate:} $y(t) = \Pr(\texttt{verdict}=\texttt{causal} \mid \text{applicability} \geq t)$
\end{itemize}

\paragraph{Elbow Rule (Threshold Selection).}
For each pair of adjacent thresholds $(t_i, t_{i+1})$ moving from stricter to looser (i.e., descending $t$), compute:
\[
\text{score}_i = \frac{\Delta x}{\lvert \Delta y \rvert} = \frac{x(t_{i+1}) - x(t_i)}{\lvert y(t_{i+1}) - y(t_i) \rvert}
\]
where $\Delta x > 0$ is the coverage gain and $\lvert \Delta y \rvert$ is the absolute drop in correct-causal rate. Select $t^* = t_{i+1}$ for the segment with maximum score. Intuition: choose the threshold where coverage gain per unit correctness loss is maximised (the ``elbow'' or diminishing-returns point).

\paragraph{Extrapolation.}
Two-stage: (1) compute pass rate and causal rate from validated sample at $t^*$; (2) project to full dataset ($n=285$) with conservative lower CI bounds. Compute projected F1 under Scenarios 1 and 2.

\paragraph{Caveats.}
Small strata (e.g., Social Norms $n=3$) have high sampling variance. Results reported with point estimates and conservative bounds.

\paragraph{Results.}
Figure~\ref{fig:rq3_human_validation_elbow}, Table~\ref{tab:rq3_human_validation_tradeoff}, Table~\ref{tab:enriched_gt_validated} (Chapter~\ref{ch:results}).

\paragraph{Reproducibility.}
\texttt{python final\_runs/RQ3\_deep\_research\_validation/scripts/enrichment\_extrapolation.py}. Analysis scripts: Appendix~\ref{sec:code_availability} (Table~\ref{tab:analysis_scripts}).

\subsubsection{Smart Trigger}
\label{sec:rq3_smart_trigger}

\paragraph{RQ \& Goal.}
Estimate a cost-efficient DR deployment policy by triggering DR only on edges flagged as uncertain by the strongest GT Lit Correctness Judge (Mechanistic prompt).

\paragraph{Source Data.}
RQ1a GT Lit Correctness outputs (Mechanistic prompt) and RQ3 Deep Research outputs. Dataset mapping: Appendix~\ref{sec:data_availability} (Table~\ref{tab:experiment_datasets}).

\paragraph{Simulation Procedure.}
Aggregate per-edge judge scores (mean over runs) from RQ1a outputs; mark edges with score $\leq 0.5$ as triggered. Evaluate the triggered subset using existing DR outputs and extrapolate expected $\Delta$F1/\$ under human-calibrated applicability thresholds (see Results Figure~\ref{fig:rq3_dr_cost_efficiency} and Methods Section~\ref{sec:rq3_threshold_selection}).

\paragraph{Simulation Parameters.}
Full configuration in Table~\ref{tab:rq3_smart_trigger_config}.

\begin{table}[H]
\centering
\caption{RQ3 Smart Trigger Configuration}
\label{tab:rq3_smart_trigger_config}
\footnotesize
\begin{tabular}{@{}ll@{}}
\toprule
\textbf{Parameter} & \textbf{Value} \\
\midrule
Trigger Threshold & mean(Aggregate Score) $\leq 0.5$ \\
Judge Used & Mechanistic Correctness (GT Lit) \\
Applicability Threshold & 0.7 \\
Aggregation & Mean score across 9 runs \\
\bottomrule
\end{tabular}
\end{table}

\textbf{Reproducibility:} Scripts: Appendix~\ref{sec:code_availability} (Tables~\ref{tab:exec_scripts},~\ref{tab:analysis_scripts}).

\endgroup
